% ID: SEL-D0017-BB-E04
% ESTADO: cotejada. Segunda pasada visual completa y borrador compilado revisados.
% Fuente: Selectividad/matematicas_aplicadas/matematicas_aplicadas_P_E.pdf
% SHA-256: bc88bfeb0f184eb0c806c4efe04e84c24f3e5fcef9caaa12e94fd48d9f67a389
% Página: 1; bloque B; ejercicio original 4.

El número de bacterias en un determinado cultivo viene dado por la
función $B(t)$, donde $t$ representa el tiempo en horas, con $0\leq t\leq 7$.
La variación instantánea en la población de bacterias en el cultivo viene
dada por la derivada de la función $B$, cuya expresión es $B'(t)=50\,000\cdot e^{2t}$.
\begin{enumerate}[label=\alph*)]
  \item \textbf{(0.75 puntos)} \textquestiondown{}Existe algún instante $t$
  en el que el número de bacterias en el cultivo comience a decrecer?
  \item \textbf{(1.5 puntos)} Obtenga la expresión de la función $B(t)$,
  sabiendo que en el instante $t=0$ el número de bacterias en el cultivo
  era de $40\,000$.
  \item \textbf{(0.25 puntos)} \textquestiondown{}Cuál es el número de
  bacterias en el cultivo a la hora y media?
\end{enumerate}
